\section{Boundaries}\label{sec:lim}

The positive results combine four ingredients: retention of grid cells by a
rounding allowance, objective-only linear noise, a choice of which
coefficients are perturbed, and exact closure certificates of specific
kinds. This section shows by explicit families how each ingredient limits
the scope. \Cref{tab:lim:summary} lists the results, including the examples
of \cref{sec:rec} that concern the same questions.

None of these results is an unconditional lower bound for all algorithms.
Each is either a limitation of a specified mechanism or a conditional
complexity implication. Some mechanism bounds hold on every draw; others
concern an expected count or an event of positive probability. Several of the instances are
easy for a different method, and the text says so. The purpose is to show
that the restrictions of the positive theorems are not artifacts of their
proofs: removing them requires a different algorithmic idea, and in some
cases it would resolve an open complexity question.

\begin{table}[tbp]
\centering
\small
\begin{tabularx}{\linewidth}{@{}>{\raggedright\arraybackslash}p{0.2\linewidth}
>{\raggedright\arraybackslash}X>{\raggedright\arraybackslash}p{0.27\linewidth}@{}}
\toprule
Result & What it limits & Kind\\
\midrule
\Cref{thm:lim:width} & Global-error retention in the sparse route: more than
$(5n/(6p))^{p/2}$ retained cells & Mechanism-specific, every draw\\
\Cref{prop:lim:local}, \cref{ex:rec:star} & Replacing the global allowance by
a bag-local one & The modified rule deletes all optimizers\\
\Cref{prop:lim:conditional} & Certified conditional values remove the
ambient dimension from the count & Sufficient interface\\
\Cref{thm:lim:ambient} & Dimension-free counts under independent uniform
noise in the low-rank route & Count surrogate, in expectation\\
\Cref{thm:lim:constraints} & Objective noise with arbitrary sparse affine
constraints & Would give $\mathrm{NP}\subseteq\mathrm{ZPP}$\\
\Cref{ex:lim:coupled}, \cref{ex:rec:feasibility} & Counting and corner bounds
when feasibility depends on the bag or core & Counting argument fails\\
\Cref{prop:lim:threshold} & Reading original threshold answers from sampled
instances & Probability bound\\
\Cref{ex:rec:fiber} & Full-dimensional convex closure under core-only noise &
Mechanism-specific, every draw\\
\Cref{prop:lim:value} & Converting value accuracy into residual distance;
Tikhonov selection & Unconditional, method-specific\\
\Cref{thm:lim:posslp} & Convex quartic residual point output under
core-only noise; treewidth two and bounded coefficients for Square Root Sum
& Would give Las Vegas algorithms for Square Root Sum and $\PosSLP$\\
\Cref{prop:lim:flow} & One-flow certificates at boundary core optima &
Mechanism-specific, probability $1/16$\\
\Cref{prop:rec:rank} & Reducing small-core recourse to a fixed quadratic
convexifier & Structural separation\\
\bottomrule
\end{tabularx}
\caption{Boundary results. ``Mechanism-specific'' results show that a stated
rule or certificate fails; they do not bound the work of other algorithms.}
\label{tab:lim:summary}
\end{table}

\subsection{The width exponent of global-error retention}\label{sec:lim:width}

The sparse theorems have bounds of the form
$C^p[c+(1+n/2)Lw_{\max}/(2\sigma)]^p\poly_d(I)$. The factor $n$ inside the
bracket comes from the rounding allowance: rounding a feasible point to the
corners of its cell changes the objective by up to $nLh^2/8$ in expectation,
and every coordinate contributes (\cref{sec:sp:width}). We show that, for
the retention rule as specified, the resulting power of $n$ is real.

\paragraph{The global-error search.}
We recall the search of \cref{sec:sp} in the special case of a continuous
box $X=[0,1]^n$, with a tree decomposition whose bags $\beta\subseteq[n]$
have at most $p$ coordinates and contain every monomial scope of $F_0$, and
with a rational $L\ge\max_i\sup_X\partial_{ii}F_0$. At level $j=0,1,\ldots$
the mesh is $h_j=2^{-j}$ and the grid is
$\mathbb G_j=(h_j\Z\cap[0,1])^n$. A cell of bag $\beta$ at level $j$ is a
product of intervals $[t_ih_j,(t_i+1)h_j]$, $i\in\beta$. At level $0$ each
bag has the single candidate cell $[0,1]^\beta$; at level $j\ge1$ the
candidate cells are the children of the cells retained at level $j-1$. The
\emph{allowed grid} $A_j$ is the set of points of $\mathbb G_j$ whose
projection onto every bag is a corner of a candidate cell of that bag. The
search computes the exact bag min-marginals
$m_\beta(v)=\min\{F_\gamma(y):y\in A_j,\ y_\beta=v\}$, updates the
incumbent value $U_j$, the least value of a feasible point found so far,
and retains a candidate cell $C$ of bag $\beta$ if and only if
\begin{equation}\label{eq:lim:retain}
 q_C-E_j\le U_j,\qquad
 q_C=\min_{v\text{ a corner of }C}m_\beta(v),\qquad
 E_j=\frac{nLh_j^2}8 .
\end{equation}
Closure is attempted on the hull $Q_j=\prod_i[a_i,b_i]$, where $[a_i,b_i]$
is the intersection, over the bags containing $i$, of the smallest interval
containing the $i$th projections of their retained cells. A coordinate is
fixed at its original lower (upper) bound only if a sound enclosure of
$\partial_iF_\gamma$ over $Q_j$ is strictly positive (negative). On the
remaining box a convexity test certifies a positive lower bound on the
Hessian of the restricted objective. If no closure succeeds by the stopping
depth $J$, the exact fallback runs. This is \cref{def:sp:algorithm} for a
continuous box. The statements below hold for every implementation of these
rules, every sound gradient enclosure, and every sound convexity test.

\begin{theorem}[Global-error retention forces a width-dependent exponent]\label{thm:lim:width}
Let $n\ge p\ge2$. Let $\mathcal G$ be the graph on $[n]$ formed by a
clique on $\{1,\ldots,p\}$ and the path $p,p+1,\ldots,n$, with edge set
$\mathcal E$. Put $\phi(t)=t^2(1-t)^2$, $\eta=1/(100n^2)$, and
\[
 F_0(x)=\sum_{i=1}^n\phi(x_i)+\eta\sum_{\{i,l\}\in\mathcal E}x_ix_l
 \qquad\text{on }X=[0,1]^n.
\]
\begin{enumerate}[label=(\alph*)]
\item\label{it:lim:width-instance} $F_0$ is a quartic whose interaction
graph is $\mathcal G$, which has treewidth $p-1$. The bags $\{1,\ldots,p\}$
and $\{i,i+1\}$, $p\le i<n$, form a tree decomposition of bag size $p$
containing every monomial scope. For every $\gamma\in\R^n$,
$\partial_{ii}F_\gamma\le2$ on $X$, so $L=2$ is valid. All widths are one,
and the base length is $I=O((n+p^2)\log n)$.
\item\label{it:lim:width-retain} Let $\gamma\in[-\frac1{100},\frac1{100}]^n$
be arbitrary, and run the global-error search on $F_\gamma$ with $L=2$. At
every level $j$ with $h_j^2\ge3p/(10n)$, every candidate cell of every bag
is retained, the hull $Q_j$ is $[0,1]^n$, no sound gradient-sign test fixes
a coordinate, and no sound convexity test succeeds.
\item\label{it:lim:width-count} Let $j_*$ be the last level with
$h_{j_*}^2\ge3p/(10n)$. If the stopping depth satisfies $J>j_*$, the search
generates more than $(5n/(6p))^{p/2}$ retained cells of the clique bag at
level $j_*$, on every such draw. The stopping depth of
\cref{def:sp:schedule} with $\sigma\le\frac1{100}$ satisfies $J>j_*$.
\item\label{it:lim:width-fpt} Consequently, for every constant $c$ there is
no function $f$ such that the work of the global-error search is at most
$f(p)I^c$ on all instances of this family.
\end{enumerate}
\end{theorem}

\begin{proof}
\ref{it:lim:width-instance}
Expanding $\phi(x_i)=x_i^2-2x_i^3+x_i^4$ shows that the only monomials
involving two variables are the products $x_ix_l$ with $\{i,l\}\in\mathcal E$,
each with coefficient $\eta\ne0$. Hence the interaction graph is $\mathcal
G$. Every tree decomposition of a graph containing a $p$-clique has a bag
containing that clique, so the treewidth is at least $p-1$. The displayed
bags, arranged as the path $\{1,\ldots,p\},\{p,p+1\},\{p+1,p+2\},\ldots$,
cover every edge and every variable, and each variable occurs in one bag or
in two adjacent bags; so they form a tree decomposition of width $p-1$. The
linear and bilinear terms have zero diagonal second derivatives, so
\[
 \partial_{ii}F_\gamma(x)=\phi''(x_i)=2-12x_i(1-x_i)\le2 .
\]
There are $O(n+p^2)$ monomials, each with indices and coefficients of
$O(\log n)$ bits, and the decomposition has $O(n+p)$ entries; this gives the
bound on $I$.

\ref{it:lim:width-retain}
Fix $\gamma$ with $\norm\gamma_\infty\le1/100$. The function
$P_\gamma(x)=\eta\sum_{\mathcal E}x_ix_l+\gamma^{\mathsf T}x$ is affine in
each coordinate separately, so moving one coordinate at a time to a better
endpoint shows that its minimum over $X$ is attained at a vertex. Since
$\phi\ge0$ and $\phi(0)=\phi(1)=0$,
\[
 f^*:=\min_XF_\gamma\ge\min_XP_\gamma=\min_{\{0,1\}^n}P_\gamma
 =\min_{\{0,1\}^n}F_\gamma\ge f^*,
\]
so $f^*$ is attained at a vertex $a$. Let $\beta$ be a bag and
$v\in[0,1]^\beta$, and put $y=(v,a_{-\beta})$. Using $\phi\le1/16$,
$\abs{\gamma_i(v_i-a_i)}\le1/100$, $\abs{y_iy_l-a_ia_l}\le1$,
$\abs{\mathcal E}\le n(n-1)/2$, and $p\ge2$,
\begin{equation}\label{eq:lim:completion}
 F_\gamma(y)-f^*
 \le\abs\beta\Bigl(\frac1{16}+\frac1{100}\Bigr)+\eta\abs{\mathcal E}
 \le\frac{29p}{400}+\frac1{200}\le\frac{3p}{40}.
\end{equation}

We show by induction that at every level $j$ with $h_j^2\ge3p/(10n)$ all
cells are candidates and all are retained. At level $0$ every cell is a
candidate. Suppose all level-$j$ cells are candidates. Then $A_j=\mathbb
G_j$. For a cell $C$ of bag $\beta$ and a corner $v$ of $C$, the point
$(v,a_{-\beta})$ lies in $\mathbb G_j$ because $a$ is a vertex, so
\eqref{eq:lim:completion} gives $q_C\le m_\beta(v)\le f^*+3p/40$. Here
$E_j=nh_j^2/4\ge3p/40$, and $U_j\ge f^*$ because $U_j$ is the value of a
feasible point. Therefore $q_C-E_j\le f^*\le U_j$, and $C$ is retained. As
every level-$j$ cell is retained, every level-$(j+1)$ cell is a candidate.

Since all cells are retained, every coordinate projection of every bag
covers $[0,1]$, and $Q_j=[0,1]^n$. Set all coordinates other than $i$ to
zero. Then $\partial_iF_\gamma=\phi'(x_i)+\gamma_i$, and
$\phi'(t)=2t(1-t)(1-2t)$ gives the values $\gamma_i+3/16>0$ at $x_i=1/4$
and $\gamma_i-3/16<0$ at $x_i=3/4$. A sound enclosure contains both values,
so it is neither strictly positive nor strictly negative, and no coordinate
is fixed. The remaining box is therefore $[0,1]^n$. At its midpoint
$c=(\frac12,\ldots,\frac12)$,
\[
 \nabla^2F_\gamma(c)=-I+\eta A_{\mathcal G},
\]
where $A_{\mathcal G}$ is the adjacency matrix. Its largest eigenvalue is at
most $-1+\eta(n-1)<0$. A function convex on a box is positive semidefinite
at interior points, so no sound convexity test can succeed.

\ref{it:lim:width-count}
Since $h_0=1$ and $3p/(10n)<1$, the level $j_*$ exists, and
$h_{j_*+1}^2=h_{j_*}^2/4<3p/(10n)$ gives $h_{j_*}^2<6p/(5n)$. The clique bag
has $h_{j_*}^{-p}=(h_{j_*}^{-2})^{p/2}>(5n/(6p))^{p/2}$ cells at level
$j_*$, all retained by \ref{it:lim:width-retain}. No closure succeeds at
levels up to $j_*$, so if $J>j_*$ the fallback does not run before level
$j_*$, and the search generates all these cells. The stopping depth of
\cref{def:sp:schedule} satisfies $h_J\le1/(4A)$ with $A\ge nL/g_0$ and
$g_0=\rho\sigma/(2W)\le\sigma/(8W)$, where $W=\sum_iw_i=n$. Here $L=2$ and
$\sigma\le1/100$, so $h_J\le\sigma/(64n^2)<\sqrt{3p/(10n)}\le h_{j_*}$, that
is, $J>j_*$.

\ref{it:lim:width-fpt}
Fix $p>2c$. For this $p$, the family has $I=O(n\log n)$ and work at least
$(5n/(6p))^{p/2}$, which exceeds $f(p)I^c$ for all sufficiently large $n$.
\end{proof}

The lower bound holds on every draw in the noise box, so it also bounds the
expected work for every law supported there, including $U_{1/100,M}$ for
every $M$. It does not contradict \cref{thm:sp:main}, whose bound has the
same form $n^{O(p)}$ for fixed numerical scales; it shows that this
exponent is necessary for the retention rule as specified. It is not a lower
bound for sparse optimization: the instances are easy, since every
optimizer is a vertex and binary dynamic programming with $O(2^p)$ states
per bag solves them. Tests on individual cells, rather than on the whole
hull, could also prune these instances much earlier.

The natural repair is to replace $E_j$ by an allowance for the bag alone,
$\abs\beta Lh_j^2/8$. \Cref{ex:rec:star} shows that this is unsound for the
unperturbed star quadratic: the error comes from minimizing over the grid
outside the bag, and it varies with the bag's own coordinates. That
deterministic example comes from an unpublished companion manuscript
\citep{companion-decomposition-aware}. The next proposition adds two points:
the failure persists on every draw of small noise, and no closure succeeds
at level $0$, so the search reaches the failing level. It is thus a failure
of the smoothed algorithm and not of one exceptional draw.

\begin{proposition}[A bag-local allowance is unsound on every draw]\label{prop:lim:local}
Let $m=32$ and let $F_0$ be the star quadratic of \cref{ex:rec:star} on
$[0,1]^{1+m}$, with center $v$, leaves $z_1,\ldots,z_m$, bags
$\{v,z_i\}$, and $L=2$. Let $\gamma\in[-\sigma,\sigma]^{1+m}$ with
$\sigma\le1/1000$, and run the search of \eqref{eq:lim:retain} with $E_j$
replaced by $\abs\beta Lh_j^2/8$ and with a stopping depth $J\ge1$, for
instance that of \cref{def:sp:schedule}. Then no closure succeeds at
level $0$, so the search reaches level $1$ ($h_1=1/2$). There every global
optimizer of $F_\gamma$ projects into the cell $C=[\frac12,1]\times[0,\frac12]$
of the bag $\{v,z_1\}$, and $C$ is deleted. The global allowance
$(1+m)Lh_1^2/8=33/16$ retains $C$.
\end{proposition}

\begin{proof}
Every $F_\gamma$ differs from $F_0$ by at most $(1+m)\sigma=33\sigma$ on the
box, and $F_0(v,z)=v^2-\frac v2\sum_iz_i+\sum_iz_i^2+\frac{15}{16}v$. The
schedule of \cref{def:sp:schedule} gives $h_J\le1/(4A)<1=h_0$, so $J\ge1$
there.

\emph{Location of the optimizers.} For fixed $v$, the coordinate $z_i$
minimizes $z_i^2-(v/2-\gamma_i)z_i$, a strictly convex function whose
minimizer on $[0,1]$ is at most $(1/2+\sigma)/2<1/2$. Let $V_\gamma(v)$ be
the minimum of $F_\gamma$ over $z$. Each inner minimum is at least
$-(v/2-\gamma_i)^2/4$, so for $v\in[0,\frac12]$
\[
 V_\gamma(v)\ge v^2+\Bigl(\frac{15}{16}-\sigma\Bigr)v-8\Bigl(\frac v2+\sigma\Bigr)^2
 =-v^2+\Bigl(\frac{15}{16}-9\sigma\Bigr)v-8\sigma^2\ge-8\sigma^2,
\]
the last step because the concave middle expression is at least its
endpoint values $-8\sigma^2$ and $7/32-9\sigma/2-8\sigma^2>0$. The feasible
point $v=1$, $z_i=1/4$ has $F_0=-1/16$, so $V_\gamma(1)\le-1/16+9\sigma<
-8\sigma^2$. Hence every optimizer has $v>1/2$ and $z_1<1/2$.

\emph{Level $0$.} At level $0$ each bag has one cell, and its value $q$ is
the minimum of $F_\gamma$ over the vertices of the box, which is at most
$F_\gamma(0)=0$. The bag-local allowance at $h_0=1$ is $1/2$, and the
incumbent is at least $\min F_\gamma\ge-1/16-33\sigma>-1/2$. So every cell
is retained, and the hull is the whole box. No closure succeeds there. Each
leaf derivative $\partial_{z_i}F_\gamma=2z_i-v/2+\gamma_i$ is negative at
$(v,z_i)=(1,0)$ and positive at $(0,1)$; the center derivative
$\partial_vF_\gamma=2v-\frac12\sum_iz_i+\frac{15}{16}+\gamma_v$ is negative at
$(v,z)=(0,\mathbf1)$ and positive at $(1,0)$. So no sound sign test fixes a
coordinate. The Hessian is the constant matrix with diagonal $2$ and
center--leaf entries $-\frac12$, and its quadratic form at
$(1,\frac14\mathbf1)$ is $2-8+4=-2$; so no sound convexity test succeeds.
Hence the search reaches level $1$, every cell of level $1$ is a candidate,
and the allowed grid is $\{0,\frac12,1\}^{1+m}$.

\emph{The level-$1$ test.} For $\gamma=0$, \cref{ex:rec:star} computes the
smallest grid min-marginal of $C$ as $23/32$; in general it is at least
$23/32-33\sigma$. The incumbent satisfies $U_1\le F_\gamma(0)=0$, because the
origin is an allowed grid point. The bag-local allowance is
$2\cdot2\cdot\frac14/8=\frac18$, and
\[
 q_C-\frac18\ge\frac{19}{32}-33\sigma>0\ge U_1 ,
\]
so $C$ is deleted. With the global allowance,
$q_C-\frac{33}{16}\le\frac{23}{32}+33\sigma-\frac{33}{16}<-1<U_1$, so $C$ is
retained.
\end{proof}

Certified values of the true conditional function remove this error, and
with it the factor $n$ in the count. The following statement makes this
precise. It is a sufficient interface, not an algorithm: computing such
values requires optimizing over all variables outside the bag. In the
notation of \cref{sec:sp:count}, let
$W_\beta(v)=V_\beta(v)+\gamma_\beta^{\mathsf T}v$ be the exact conditional
value of a bag tuple, a minimum over the original outside domain. Parts~(a)
and~(b) below restate deterministic guarantees of the companion's certified
bag filter \citep{companion-decomposition-aware}; part~(c) is the
finite-law count of \cref{prop:sp:conditional}(c).

\begin{corollary}[Certified conditional values remove the dimension factor]\label{prop:lim:conditional}
Let the modified search and certified conditional oracle of
\cref{prop:sp:conditional} be given, with an error allowance $ae_j$,
$a\ge0$, and $e_j=pLh_j^2/8$. Then:
\begin{enumerate}[label=(\alph*)]
\item\label{it:lim:cond-sound} every global optimizer is retained on every
 draw at every level;
\item\label{it:lim:cond-witness} the incumbent satisfies
 $U_j\le f^*+(1+a)e_j$, and every retained bag cell has a corner with true
 conditional value at most $f^*+2(1+a)e_j$;
\item\label{it:lim:cond-count} for $j\le J$ and $M\ge2^J$, the expected
 number of such tuples in the complete bag grid is bounded by
 \eqref{eq:sp:conditional}, with error parameter $a$.
\end{enumerate}
Thus, for fixed $a$, the count is FPT jointly in $p$ and $Lw_{\max}/\sigma$.
\end{corollary}

\begin{proof}
These are the retention, witness, and counting conclusions of
\cref{prop:sp:conditional}. Its incumbent includes the current oracle
answers before pruning, and its error includes the entire outside problem.
\end{proof}

With $a$ fixed, the per-bag factor depends on the bag size $p$ but not on
the dimension $n$, so the count would be fixed-parameter in $p$ and the
numerical ratio. What is missing for a general sparse algorithm is the
oracle: the grid min-marginals of \cref{sec:sp} do not satisfy the interface
(\cref{prop:lim:local}), and an exact conditional value can require solving
almost the whole instance. The recourse route of \cref{sec:rec} realizes the
interface for one search block, a small core, by assuming a tractable
residual problem; its count is the case of a single bag of size $k$. For
general sparse instances no such oracle is available, and whether another
method achieves a bound FPT jointly in width and numerical ratios under linear noise
remains open.

\subsection{Ambient uniform noise and dimension-free counts}\label{sec:lim:ambient}

The low-rank route has bounds FPT jointly in structure and numerical ratios under aligned noise
(\cref{thm:qp:aligned}) and under Gaussian-like noise on all coefficients
(\cref{thm:qp:gauss}), whereas under independent uniform noise on all
coefficients its count contains powers of the dimension
(\cref{thm:qp:uniform}). The reason is that the factor component and the
residual component of a uniform cube vector are dependent: conditionally on
the residual component, the factor component can be concentrated much more
than its variance suggests. For a Gaussian vector the two components are
independent. \Cref{prop:qp:ambient-barrier} states that the dimension
dependence of the uniform count is real for the counted quantities, even for
perfectly conditioned factors and fixed projected widths. We prove it here,
in a slightly more general form.

We use the notation of \cref{sec:qp:aux} with $\Lambda=\alpha I$. For a
factor $T\in\Q^{k\times N}$ with $TT^{\mathsf T}=I_k$, write
$\gamma=T^{\mathsf T}\xi+\zeta$ with $\xi=T\gamma$ and $T\zeta=0$, and
\begin{equation}\label{eq:lim:aux}
 V(a)=\min_{x\in P}\Bigl[F_0(x)+\zeta^{\mathsf T}x
 +\frac\alpha2\norm{a-Tx}^2\Bigr]+\xi^{\mathsf T}a,\qquad a\in\R^k .
\end{equation}
The search of \cref{thm:qp:uniform} refines dyadic meshes on a box that
contains the projected ranges enlarged by the largest possible factor
noise. When all its widths are equal, the allowance at mesh $h$ is
$E_h=\alpha kh^2/8$, and the count of \cref{thm:qp:uniform} bounds the
expected number of grid nodes $v$ satisfying the local event
$V(v)\le V(v\pm he_l)+2E_h$ for every coordinate $l$ that is interior at
$v$.

\begin{theorem}[Uniform ambient counts depend on the dimension]\label{thm:lim:ambient}
Let $k\ge1$, let $m\ge256$ be even, let $\alpha\in\{1,m\}$, and put
$N=km^2$. There are a symmetric rational matrix $A$, a rational factor $T$
with $TT^{\mathsf T}=I_k$ and $A+\alpha T^{\mathsf T}T\succeq0$, and a bounded
rational polytope $P\subset\R^N$ with $O(N)$ inequalities, all of encoding
length polynomial in $N$, with the following properties for
$F_0(x)=\frac12x^{\mathsf T}Ax$.
\begin{enumerate}[label=(\alph*)]
\item\label{it:lim:amb-instance} $A$ has exactly $k$ negative eigenvalues,
each equal to $-\nu$ with $\nu=\alpha(\sqrt{17/8}-1)$, so $\alpha<4\nu$, and
each coordinate of $Tx$ ranges over an interval of width $8$ as $x$ ranges
over $P$.
\item\label{it:lim:amb-count} Let the coordinates of $\gamma$ be
independent, either uniform on $[-1,1]$ or with law $U_{1,M}$ for a power of
two $M\ge32$. For every grid of spacing $h$ on an interval containing
$[-\frac32,\frac32]$, with product grid $\mathbb G_h$ in $\R^k$:
\begin{enumerate}[label=(\roman*)]
\item if $8/\sqrt{\alpha m}\le h\le16/\sqrt{\alpha m}$, then
$\E\,\#\{v\in\mathbb G_h:V(v)\le\min V+2E_h\}\ge(\sqrt{\alpha m}/64)^k$;
\item if $16/(\alpha m)\le h\le32/(\alpha m)$, then the expected number of
nodes $v\in\mathbb G_h\cap[-1,1]^k$ with $V(v)\le V(v\pm he_l)+2E_h$ for
every $l$ and both signs is at least $(\alpha m/128)^k$.
\end{enumerate}
The dyadic meshes of the search of \cref{thm:qp:uniform} include values of
$h$ in both ranges.
\item\label{it:lim:amb-fpt} For $\alpha=1$ these bounds are
$64^{-k}(N/k)^{k/4}$ and $128^{-k}(N/k)^{k/2}$. Consequently no bound of the
form $f(k,R)\,I^c$, with an absolute constant $c$, holds for either expected
count, where $R$ records the frame $TT^{\mathsf T}$, the projected widths,
the ratio $\alpha/\nu$, and the noise scale, all of which are fixed here.
\end{enumerate}
\end{theorem}

The proof is in \cref{app:lim:ambient}. Each block of the construction has a
negative direction coupled to a thin coordinate, and a simplex of total mass
$4m$ whose two halves carry the signs of the factor row. With constant
probability the smallest noise coefficients of the two halves are both
close to $-1$; then minimization over the simplex cancels the factor noise up
to $O(1/m)$, the auxiliary objective is flat to within $O(1/m)$ on a long
interval, and many grid nodes pass every comparison.

The theorem bounds the count surrogates used in the analysis, not the work
of the algorithm: the closure step could resolve these instances before the
counted nodes are visited. It also leaves room for bounds in terms of the
ambient diameter, which grows like $m$ here. What it excludes is a proof of
a fixed-parameter bound for uniform ambient noise that bounds these counts
in terms of projected widths and conditioning alone. \Cref{thm:qp:gauss}
avoids the obstruction because the factor and residual components of a
Gaussian vector are independent.

\subsection{Constraints and decision thresholds}\label{sec:lim:constraints}

The sparse results allow mixed boxes, products of simplices, order
constraints, and graph equalities (\cref{sec:con}). The next result shows
that objective noise cannot handle arbitrary sparse affine constraints:
their feasibility problem can encode an NP-complete question that no small
objective perturbation changes.

\begin{theorem}[Objective noise does not remove sparse feasibility hardness]\label{thm:lim:constraints}
There is a polynomial-time map from instances $(a_1,\ldots,a_n;B)$ of
\textup{\textsc{Subset Sum}}, with positive integers, to the following instances.
Put $A=\sum_ia_i+B$. The variables are $t,z_1,\ldots,z_n,w_1,\ldots,w_n\in
\{0,1\}$ and $s_1,\ldots,s_n\in[0,1]$; with $s_0=0$ the constraints are
\[
 z_i+w_i=t,\qquad s_i=s_{i-1}+\frac{a_i}Az_i\quad(i=1,\ldots,n),\qquad
 s_n=\frac BAt,
\]
and the objective is $F_0=-t$. Every constraint has at most three variables
and coefficients of absolute value at most one, the constraint and objective
scopes have a tree decomposition of bag size four, and the all-zero point is
feasible. For every $\gamma\in\R^{3n+1}$ with
$\norm\gamma_\infty\le\sigma:=1/(10(3n+2))$, every global optimizer of
$F_\gamma=-t+\gamma^{\mathsf T}x$ has $t=1$ if the instance is a
yes-instance and $t=0$ otherwise.

Consequently, suppose an exact smoothed algorithm handles mixed boxes with
additional affine equality constraints whose scopes fit a supplied tree
decomposition of bag size four, under ambient noise drawn from a
base-chosen finite rational law supported in $[-\sigma,\sigma]$ with
polynomially many sampling bits. If its expected work is polynomial in $I$
and $1/\sigma$, then \textup{\textsc{Subset Sum}} has a Las Vegas algorithm with
expected polynomial running time, and $\mathrm{NP}\subseteq\mathrm{ZPP}$.
\end{theorem}

\begin{proof}
If $t=0$, then $z_i+w_i=0$ forces $z_i=w_i=0$, the recurrence forces
$s_i=0$, and $s_n=0$ holds; so the all-zero point is the only feasible
point with $t=0$. If $t=1$, then $w_i=1-z_i$ and
$s_i=A^{-1}\sum_{l\le i}a_lz_l\in[0,1]$, and $s_n=B/A$ holds exactly when
$\sum_la_lz_l=B$. So a feasible point with $t=1$ exists exactly for
yes-instances.

The bags $P_i=\{t,s_{i-1},s_i,z_i\}$ (with $s_0$ omitted) form a path, and
attaching $\{t,z_i,w_i\}$ to $P_i$ gives a tree. Every constraint lies in a
bag, $t$ lies in every bag, and every other variable lies in one bag or in
two adjacent bags; so this is a tree decomposition of bag size four.

There are $3n+1$ variables in $[0,1]$, so
$\abs{\gamma^{\mathsf T}x}\le(3n+1)\sigma<1/10$. The zero point has value
$0$, and every feasible point with $t=1$ has value below $-9/10$. Hence the
optimizers have $t=1$ exactly on yes-instances.

For the consequence, construct the instance, compute and sample the law,
run the algorithm, and read $t$ from its exact output. The answer is
correct on every draw, and the expected running time is polynomial because
$1/\sigma=O(n)$ and the instance has length polynomial in the
\textsc{Subset Sum} input.
\end{proof}

The objective is linear, so any positive $L$ is a valid curvature bound, all
widths are one, and $1/\sigma=O(n)$. A constrained analogue of the sparse
bounds with these parameters would therefore be polynomial on this family.
Bounded coefficient magnitudes do not bound denominators, Hoffman constants,
or the constants of feasible repair, and the reduction makes no claim about
them. A constrained theorem can avoid the barrier by assuming structure
that preserves feasibility under rounding, as the simplex, order, and graph
results of \cref{sec:con} do, or by charging parameters that are large on
this family.

Constraints also affect the counting argument itself. The box count uses
the fact that a conditional value, obtained by minimizing over the variables
outside a bag or a core, has bounded upper curvature in each coordinate.
This requires the outside feasible set to be independent of the bag or core
values. \Cref{ex:rec:feasibility} shows that otherwise the corrected-corner
lower bound itself fails. The following example shows that the expected
count fails as well, even for a single totally unimodular row.

\begin{example}[Core-dependent feasibility defeats the coordinatewise count]\label{ex:lim:coupled}
Let $v_1,v_2,y\in[0,1]$ with the single constraint $y\ge v_1-v_2$, which
together with the bounds is totally unimodular. Minimize
$2y+\gamma_1v_1+\gamma_2v_2+\gamma_yy$, with independent coefficients either
uniform on $[-\frac12,\frac12]$ or with law $U_{1/2,M}$. Regard $(v_1,v_2)$
as a core, or as a bag, and $y$ as the residual variable. For fixed
$\gamma_y$ the conditional value is
\[
 V(v)=\kappa\max\{0,v_1-v_2\},\qquad\kappa=2+\gamma_y\in[\tfrac32,\tfrac52],
\]
which has a convex kink along the diagonal. Let $h=1/r$ and consider the
diagonal grid nodes $(kh,kh)$, $1\le k\le r-1$. Their four coordinate
neighbors increase $V+\gamma_1v_1+\gamma_2v_2$ by
\[
 h(\kappa+\gamma_1),\qquad-h\gamma_1,\qquad h\gamma_2,\qquad h(\kappa-\gamma_2).
\]
On the event $\{\gamma_1<0<\gamma_2\}$, which has probability $1/4$ under
both laws because $0$ is not an atom of $U_{1/2,M}$, all four increments are
positive. So every one of the $r-1$ diagonal nodes passes every
coordinate-neighbor comparison, even with tolerance zero, and the expected
number of passing nodes is at least $(r-1)/4$, unbounded as $h\to0$. Under a
coordinate upper-curvature bound $L$, the two increments in one coordinate
sum to at most $Lh^2$, which confines the coefficient to an interval of
length $O(Lh)$; here they sum to $\kappa h$.
\end{example}

The example does not bound the number of truly near-optimal nodes: along the
diagonal the objective is $kh(\gamma_1+\gamma_2)$, and comparisons in the
diagonal direction remove the spurious candidates. The constrained sparse
results therefore use feasible rounding and comparisons along the faces of
their polytopes, and the recourse results require a residual feasible set
that does not depend on the core.

Finally, the answer to a threshold question about the base instance cannot
in general be read off from the exact optimum of a sampled instance when the
threshold gap is smaller than the noise. Worst-case hardness results for
structured nonconvex optimization, such as strong NP-hardness of
box-constrained quadratic programming at treewidth
two~\citep{pia2026-treewidth-and-the-complexity-of}, concern such threshold questions, and bounded
coefficients do not by themselves give a lower bound on the distance of a
no-instance's minimum from the threshold.

\begin{proposition}[Small gaps do not survive the noise]\label{prop:lim:threshold}
Let $\Psi$ be continuous on a compact set $X$, let $\delta\ge0$, and let
$y\in X$ satisfy $\Psi(y)\le\delta$ and $\abs{y_i}\ge1$ for some coordinate
$i$. Let
$\gamma$ have independent coordinates with laws symmetric about zero, and
let $\gamma_i$ be uniform on $[-\sigma,\sigma]$. Then
\[
 \Pr\Bigl\{\min_X(\Psi+\gamma^{\mathsf T}x)<0\Bigr\}\ge\frac12-\frac\delta{2\sigma}.
\]
If instead $\gamma_i$ has law $U_{\sigma,M}$, the bound holds with the
additional term $-1/(2M)$.
\end{proposition}

\begin{proof}
Put $Z=\gamma^{\mathsf T}y$. It is a sum of independent symmetric terms, so
$Z$ is symmetric and $\Pr\{Z<-\delta\}=\frac12(1-\Pr\{\abs Z\le\delta\})$.
Conditionally on $(\gamma_l)_{l\ne i}$, the event $\abs Z\le\delta$ confines
$\gamma_i$ to an interval of length $2\delta/\abs{y_i}\le2\delta$, which has
probability at most $\delta/\sigma$, or at most $\delta/\sigma+1/M$ by
\eqref{eq:model:interval}. On the event $Z<-\delta$,
$\min_X(\Psi+\gamma^{\mathsf T}x)\le\Psi(y)+Z<0$.
\end{proof}

If $\Psi\ge0$ and $0<\min_X\Psi\le\delta\ll\sigma$, the base instance is a
no-instance of the question ``is $\min\Psi\le0$?'', while the sampled
instance has a negative minimum with probability close to $1/2$. Reading
the original answer directly from the sampled optimum therefore fails on
such instances, and expected polynomial bounds for sampled instances are
compatible with strong NP-hardness of the corresponding threshold
questions. These guarantees do not by themselves provide a polynomial-time
algorithm for the original threshold questions. The proposition does not
exclude other uses of a smoothed solver, for instance reductions with a
robust gap or repeated queries at chosen noise scales.

\begin{remark}[The treewidth-two reduction]\label{rem:lim:dk}
The strong NP-hardness proof of \citet{pia2026-treewidth-and-the-complexity-of}
reduces \textsc{Subset Sum} to the question whether a nonnegative box
quadratic $\Psi$ of treewidth two with bounded integral coefficients has
minimum zero. In that construction $\Psi$ combines nonnegative binary
penalties $x_i(1-x_i)$ with squared residuals for bit copies, bit-serial and
partial-sum recurrences, and a final square that compares the selected
partial sum with the target. Take
$r\ge1$, $A=2^r$, items $A$ and $A+2$, and target $A+1$, a no-instance. Selecting
the first item and following every recurrence exactly gives a point of the
unit box at which every term vanishes except the final square, which equals
$2^{-2r-4}$, and the indicator of the selected item equals one.
\Cref{prop:lim:threshold} with $\delta=2^{-2r-4}$ shows that, under
independent noise of half-width $\sigma$ with law $U_{\sigma,M}$, the sampled
minimum is negative with probability at least
$\frac12-2^{-2r-5}/\sigma-1/(2M)$, although the base instance is a
no-instance. The number of variables of the construction is linear in $r$,
so the gap is exponentially small in the instance size.
\end{remark}

\subsection{Core-only noise and residual outputs}\label{sec:lim:points}

Core-only noise perturbs the few coordinates that carry the nonconvexity
and leaves the residual costs unchanged. Ties and flat directions in the
residual variables then persist on every draw. The continuous core-only
theorem (\cref{thm:rec:core-only}) assumes a certified uniform
strong-convexity modulus of the residual problem, and the integer core-only
results of \cref{sec:int} use a certificate that represents all tied
optimal flows at once. These are sufficient structures, rather than
necessary conditions for every residual problem. The exact-arithmetic
companion~\citep{companion-exact-arithmetic} also gives full
selected-optimizer oracles under core-only noise for residual-convex cubics,
without a residual modulus. Its supplied-convexifier theorem also gives
point oracles when the transformed objective is cubic on the feasible
polytope or is a fixed-degree polynomial convex on all of $\R^n$.
The general arithmetic obstruction below begins at degree
four. The examples explain the limits of the certificates used here.

\Cref{ex:rec:fiber} concerns the closure step used under ambient noise: a
full-dimensional box around the optimal set on which the objective is
convex. Under core-only noise such a box need not exist, even when the core
has projected growth one, the residual problem is convex or strictly convex,
and the optimizer is unique and interior. The example has a short exact
optimality certificate, so it limits one certificate, not exact
optimization. Two further results concern the output itself. The first
separates value accuracy from distance to an optimizer. It shows that no
value tolerance with polynomially many bits localizes the residual
coordinates of a convex quartic with constant-size coefficients, and that a
noisy core does not change this.

\begin{proposition}[Value accuracy does not localize residual points]\label{prop:lim:value}
For $n\ge1$ and $x_0=\frac12$, let
\[
 G_n(x)=\sum_{i=1}^n\bigl(8x_i^2-x_{i-1}^2x_i\bigr),\qquad
 \Phi_\gamma(v,x,y)=\Bigl(v-\frac12\Bigr)^2+\gamma v+G_n(x)+x_n^2(y-1)^2
\]
on $[0,1]\times[0,1]^n\times[0,1]$, with $\abs\gamma\le\frac12$.
\begin{enumerate}[label=(\alph*)]
\item\label{it:lim:val-convex} $\Phi_\gamma$ is a convex quartic with
$O(n)$ monomials and constant-size coefficients other than $\gamma$, and
$\nabla^2G_n\succeq10I$ on $[0,1]^n$.
\item\label{it:lim:val-opt} $\Phi_\gamma$ has the unique minimizer
$(v_\gamma,x^*,1)$, where $v_\gamma=(1-\gamma)/2$, $x^*$ is the interior
minimizer of $G_n$, and $0<x_n^*\le2^{-2^{n+1}}$.
\item\label{it:lim:val-gap} The feasible point $(v_\gamma,x^*,0)$ is at
distance one from the minimizer, and its objective gap is
$(x_n^*)^2\le2^{-2^{n+2}}$.
\item\label{it:lim:val-tikhonov} For $\lambda>0$ let $(x_\lambda,y_\lambda)$
minimize $G_n(x)+x_n^2(y-1)^2+\lambda(\norm x^2+y^2)$ on $[0,1]^{n+1}$. If
$y_\lambda\ge\frac12$, then $\lambda\le2^{-2^{n+2}}$.
\end{enumerate}
\end{proposition}

\begin{proof}
\ref{it:lim:val-convex}
The Hessian of $G_n$ is tridiagonal, with diagonal entries $16-2x_{i+1}$
for $i<n$ and $16$ for $i=n$, and off-diagonal entries $-2x_i$. On the box
each diagonal entry is at least $14$ and each row has off-diagonal absolute
sum at most $4$, so $\nabla^2G_n\succeq10I$ by Gershgorin's theorem. For a
direction $(p,q)$ in $(x,y)$ and $u=1-y\in[0,1]$, direct differentiation
gives
\[
 (p,q)^{\mathsf T}\nabla^2\bigl[x_n^2(y-1)^2\bigr](p,q)
 =2(x_nq-2up_n)^2-6u^2p_n^2,
\]
so the Hessian of $G_n+x_n^2(y-1)^2$ has quadratic form at least
$10\norm p^2-6p_n^2\ge0$. The core term is convex and separate.

\ref{it:lim:val-opt}
The derivative $\partial_iG_n=16x_i-x_{i-1}^2-2x_ix_{i+1}$ (the last term
absent for $i=n$) equals $-x_{i-1}^2$ at $x_i=0$ and is at least $13$ at
$x_i=1$. Since $x_0>0$, induction on $i$ shows that the strongly convex
function $G_n$ has its minimizer $x^*$ in the interior, with
$x_i^*=(x_{i-1}^*)^2/(16-2x_{i+1}^*)\le(x_{i-1}^*)^2/14$, where
$x_{n+1}^*:=0$. Writing $x_i^*=14e_i$ gives $e_i\le e_{i-1}^2$ and
$e_0=1/28$, so $x_n^*\le14\cdot28^{-2^n}\le4^{-2^n}$. Since
$\Phi_\gamma\ge(v-\frac12)^2+\gamma v+G_n(x)$, with equality exactly when
$x_n^2(y-1)^2=0$, and $x_n^*>0$, the unique minimizer is
$(v_\gamma,x^*,1)$.

\ref{it:lim:val-gap} is immediate from \ref{it:lim:val-opt}.

\ref{it:lim:val-tikhonov}
The regularized objective is strongly convex, so its minimizer is unique.
For fixed $x$, minimizing $x_n^2(y-1)^2+\lambda y^2$ over $y$ gives
$y_\lambda=x_{n,\lambda}^2/(x_{n,\lambda}^2+\lambda)$. The boundary argument
of \ref{it:lim:val-opt} applies unchanged to the regularized problem, and
its stationarity equations give
$x_{i,\lambda}=x_{i-1,\lambda}^2/(16+2\lambda-2x_{i+1,\lambda})$ for $i<n$
and $x_{n,\lambda}=x_{n-1,\lambda}^2/(16+2\lambda+2(1-y_\lambda)^2)$. Hence
$x_{n,\lambda}\le4^{-2^n}$ as before. If $y_\lambda\ge\frac12$, then
$\lambda\le x_{n,\lambda}^2\le2^{-2^{n+2}}$.
\end{proof}

The input has length $O(n\log n)$, so part~\ref{it:lim:val-gap} says that a
value accuracy with fewer than $2^{n+2}$ bits, which is superpolynomial in
$I$, certifies no residual accuracy below one, on every draw.
Part~\ref{it:lim:val-tikhonov} shows that isotropic Tikhonov selection, the
usual way to pick a point from a convex optimal fiber, needs a
regularization parameter of exponentially many bits here. The instance is
nonetheless easy: setting $y=1$ leaves a uniformly strongly convex problem.
So the proposition limits methods that convert values into points, not
residual optimization in general.

The next result concerns every method. It encodes two arithmetic decision
problems in the residual optimizer of a convex quartic. \emph{Square Root
Sum} asks, for positive integers $a_1,\ldots,a_n$ and an integer $B$, whether
$\sum_i\sqrt{a_i}\le B$. A \emph{straight-line program} computes a sequence
of integers, starting from the constants $0$ and $1$, each a sum, difference,
or product of two earlier ones; its output is the last integer $A$, and
$\PosSLP$ asks whether $A>0$. Both problems lie in the counting hierarchy,
and Square Root Sum is decidable in polynomial time with a $\PosSLP$
oracle~\citep{allender2009-on-the-complexity-of-numerical}; whether either
is in $\mathrm P$ is open.

\begin{theorem}[Residual points under core-only noise]\label{thm:lim:posslp}
\leavevmode
\begin{enumerate}[label=(\alph*)]
\item\label{it:lim:pt-srs} There is a polynomial-time algorithm that maps
every Square Root Sum instance either to its answer or to an explicit
rational quartic $\Phi$ on a unit box $Y$, convex on $Y$, with a unique
minimizer $\hat x$ whose designated coordinate $\hat x_y$ is $1$ if
$\sum_i\sqrt{a_i}>B$ and $0$ if $\sum_i\sqrt{a_i}<B$; equality does not
occur in the produced instances. The interaction graph of $\Phi$ has
treewidth at most two, every nonconstant coefficient has absolute value
bounded by an absolute constant, and $\partial_{ii}\Phi\le20$ on $Y$.
\item\label{it:lim:pt-posslp} There is a polynomial-time algorithm that maps
every straight-line program with output $A$ to an explicit rational quartic
$\Phi$ on a rational box $Y$, convex on $Y$, with a unique minimizer $\hat x$
whose designated coordinate is $1$ if $A>0$ and $0$ if $A\le0$.
\item\label{it:lim:pt-core} For $\Phi$ from (a) or (b), consider
\[
 F_\gamma(v,x)=\Bigl(v-\frac12\Bigr)^2+\gamma v+\Phi(x)\qquad\text{on }[0,1]\times Y,
\]
with one noisy core coordinate $v$, no residual noise, and
$\abs\gamma\le\frac12$. The core optimizer is $v_\gamma=(1-\gamma)/2$, the
projected growth is one, $\partial_{vv}F_\gamma=2$, the residual objective
is convex, and the unique optimizer is $(v_\gamma,\hat x)$ on every draw.
Consequently, if an exact smoothed algorithm for this class, under a
base-chosen finite rational core-noise law with polynomially many sampling
bits, returns an output from which a point within distance $1/4$ of the
optimizer can be computed with expected polynomial total work, then Square
Root Sum and $\PosSLP$ have Las Vegas algorithms with expected polynomial
running time; for Square Root Sum the instances of (a), with treewidth two
and bounded coefficients, suffice. If a deterministic algorithm returns, in
polynomial time, a point within distance $1/4$ of the minimizer of every
explicit rational quartic that is convex on a rational box and has a unique
minimizer, then Square Root Sum and $\PosSLP$ are in $\mathrm P$.
\end{enumerate}
\end{theorem}

The proof is in \cref{app:lim:posslp}. All three parts, including the
core-only-noise and deterministic consequences, also appear in the
unpublished companion manuscript~\citep{companion-exact-arithmetic}; the
constructions are reproduced here to keep the paper self-contained. The two parts have different structure:
part (a) has treewidth two and bounded coefficients, while the construction
of part (b) bounds neither its treewidth nor its coefficient magnitudes. The
theorem is a conditional implication, not an NP-hardness result. Values
remain easy: convex optimization approximates $\min\Phi$ in polynomial time,
and the program $\min_Y\Phi$ is itself a compact exact description of
$\hat x$. What the theorem limits is the efficient evaluation of residual
coordinates. The residual problem has no positive uniform strong-convexity
modulus on its box: along the designated coordinate its curvature vanishes
on the face where both activation coordinates are zero, and at $\hat x$ it
can be extremely small. This is consistent with the certified modulus
assumed in \cref{thm:rec:core-only}. Part (a) is also relevant to the sparse
route: its residual problems have treewidth two and bounded coefficients,
so sparsity of the residual problem does not by itself make residual points
efficiently computable under core-only noise. Noise on the residual
coordinates is what supplies point growth in the ambient theorems.

For integer recourse under core-only noise, the analogue of a convex patch
is a certificate that one residual label, for instance one flow, is optimal
throughout a small core box. Boundary core optima defeat this certificate on
an event whose probability does not shrink as the noise grid is refined.

\begin{proposition}[One-flow certificates fail at boundary core optima]\label{prop:lim:flow}
Let $v\in[0,1]^2$, let one integral unit be routed through two parallel
arcs, so that the residual labels are $Y=\{(1,0),(0,1)\}$, and let
\[
 F_\gamma(v,z)=\tfrac12\norm v^2+\gamma^{\mathsf T}v+v_1z_1+v_2z_2 ,
\]
with core-only noise $\gamma_1,\gamma_2$ independent with law $U_{1,M}$,
$M\ge4$. Then $L=1$ is valid and the event
$\mathcal E=\{\gamma_1\ge\frac12,\gamma_2\ge\frac12\}$ has probability
exactly $1/16$. On $\mathcal E$:
\begin{enumerate}[label=(\alph*)]
\item\label{it:lim:flow-growth} the projected value
$V_\gamma(v)=\frac12\norm v^2+\gamma^{\mathsf T}v+\min\{v_1,v_2\}$
satisfies $V_\gamma(v)\ge(\frac12+\min_i\gamma_i)\norm v^2$, so $v=0$ is
the unique optimal core, with growth at least one, and both labels are
optimal there;
\item\label{it:lim:flow-box} on every box $[0,a]\times[0,b]$ with $a,b>0$ no
single label is optimal throughout, so no certificate of optimality of a
single label on such a box exists;
\item\label{it:lim:flow-search} the corrected-grid core search with
allowance $kLh^2/8=h^2/4$ retains the cell $[0,h]^2$ at every level.
\end{enumerate}
If the two arcs are followed by $m-1$ further stages of two parallel arcs
with zero cost, there are $2^m$ feasible flows and the same conclusions hold.
A strategy that requires a one-label certificate on the retained hull and
otherwise enumerates all flows then has expected work at least $2^m/16$,
while the input length is $O(m\log m)$ plus the sampling bits.
\end{proposition}

The proof is in \cref{app:lim:flow}. The instance is easy: on $\mathcal E$
the bound in \ref{it:lim:flow-growth} certifies the origin and any flow. The
integer core-only theorem of \cref{sec:int} avoids the obstruction with a
certificate that represents every optimal flow at once and fixes the core
face instead of one label.

\subsection{Relations between the structural routes}\label{sec:lim:rank}

The low-negative-inertia route uses a quadratic correction
$\alpha\norm{Tx}^2/2$ of small rank $k$ that makes the objective convex. The
recourse route uses a small core whose fixing makes the residual problem
tractable. \Cref{prop:rec:rank} gives a family with a core of size one and
convex polynomial recourse for which every fixed positive semidefinite
quadratic correction that makes the objective convex has rank at least the
residual dimension. \Cref{thm:rec:poly} applies to it with the core
curvature as numerical parameter, independently of the coupling scale, while
a quadratic convexification would need rank equal to the residual
dimension. The two routes therefore measure different structure. The separation concerns the
broader class with merely convex recourse. A uniformly strongly convex
residual problem, as in \cref{thm:rec:core-only}, does admit a quadratic
convexifier supported on the core, of rank at most $k$, by a Schur-complement
bound; for that class the contribution of the recourse theorem is that its
numerical factor keeps the original core curvature $L/\sigma$ rather than the
scale of such a convexifier. If the full Hessian blocks are bounded in
operator norm by $M$ and the residual modulus is $\mu$, a sufficient
convexifier scale is of order $M+M^2/\mu$. The direct theorem retains the
original $L/\sigma$ in its count. Neither statement is a hardness result.

The sparse and recourse routes are related in the opposite way. A bag of the
sparse route plays the role of a core, and the sparse count contains the
factor $n$ only because it lacks exact conditional values for the variables
outside the bag (\cref{prop:lim:local,prop:lim:conditional}). Where such
values are available, as in the recourse theorems, the factor $n$
disappears. The strong-noise
route of \cref{sec:int} uses no structural parameter at all, and its noise
regime is different from the other three.
